% problem-000214 7 铯碳酸盐分解热力学
\section*{7. 铯碳酸盐分解热力学}
\textit{下列为分问。}

\subsection*{7-1}
分别画出三氟甲氧基苯和2-甲氧基吡啶的优势构象。

\subsection*{7-2}
试⽐较以下两组⾃由基断裂化反应的相对速率（断裂的键已⽤箭头标出），并简要解释你的判断。

\subsection*{7-2}
-1

A

B

\subsection*{7-2}
-2

（图：）

（图：）

B

\subsection*{7-3}
给以下三个化合物的电离势从⾼到低进⾏排序。

（图：）



（图：）


（图：）



\subsection*{7-4}
在1H-NMR表征中，H核受屏蔽作⽤影响展现不同的化学位移。屏蔽作⽤随电⼦密度增⼤⽽增⼤，对应H的化学位移 (ppm) 数值则减⼩。

\subsection*{7-4}
-1 将以下两种化合物中⼼C原⼦所连接H的化学位移从⼤到⼩进⾏排列。

（图：）



（图：）



\subsection*{7-4}
-2 画出以下两个化合物的⽴体结构，并对箭头所指H的化学位移从⼤到⼩进⾏排列。

（图：）



（图：）



\subsection*{7-5}
在有机合成中，常需要使⽤⽆⽔溶剂，如四氢呋喃常⽤与⾦属钠混合回流，确定⽆⽔后重蒸即可。以下可⽤于此四氢呋喃/钠体系达致⽆⽔的指⽰剂的有：( )a) ⼆苯酮； b) 苯酚； c) 苯甲醛； d) 苯⼄酮

\subsection*{7-6}
在有机化合物的提纯⽅法中，蒸馏是常⻅的⽅法。以下可利⽤⽔蒸⽓蒸馏法提纯的化合物有：( )a) ⼄酸⼄酯； b) 硝基苯； c) 苯甲酸⼄酯； d) 苯胺第8题（22分, 13%）杂原⼦参与的ene反应

烯丙基⼆氮烯重排反应 (Allylic Diazene Rearrangement, ADR) 可认为是 ene 反应的逆反应，在脱去氮⽓的同时，发⽣双键的位移。其反应式如下所⽰：

（图：）

\subsection*{8-1}
若底物结构中有⼿性碳原⼦时，反应具有⽴体选择性。画出中间体A的⽴体结构式，判断主要产物B中⼿性碳原⼦C1和C2的绝对构型及双键的构型。

（图：）

\subsection*{8-2}
Retro-ene 反应和 Diels-Alder 反应相结合，可以构筑环状 1,4-⼿性碳原⼦。画出反应主要产物 A、B 以及C的结构，判断D中两个甲基的相互关系，以及硫元素最后的存在形式。

（图：）
8-3 依据以上信息，完成以下反应：

\subsection*{8-3}
-1

（图：）
